\chapter{Problèmes}
\section{Synthèse}
\begin{exo}
\index{Fibonacci!Suite de}
	On d\'efinit la suite de Fibonacci $(F_n)_{n\in\mathbb{N}}$ par :
	\[
		\begin{multiform}
			F_0 = 0 \\
			F_1 = 1 \\
			F_{n+2} = F_{n+1} + F_{n} .
		\end{multiform}
	\]
	Montrer que, pour tous entiers naturels $a$ et $b$ tels que $a \geq b$, on a
	$\pgcd{a}{b} = \pgcd{b}{a-b}$.
	En d\'eduire que deux nombres de Fibonacci cons\'ecutifs, c'est-\`a-dire
	$F_n$ et $F_{n+1}$, sont premiers entre eux.
\end{exo}

\begin{exo}
	On lance un d\'e dod\'eca\`edrique (\`a 12 faces) r\'egulier non-truqu\'e dont 
	les faces sont num\'erot\'ees de z\'ero \`a onze.
	On s'int\'eresse au nombre \'ecrit sur la face sup\'erieure du d\'e.
	Les nombres sur ses faces sont \'ecrits en base trois.
	\begin{enumerate}
		\item 
			\begin{enumerate}
				\item Montrer qu'un nombre \'ecrit en base trois est un multiple de $3$ si et seulement s'il se termine par $0$.
				\item Calculer alors la probabilit\'e que le nombre obtenu par le lancer du d\'e se termine par $0$.
			\end{enumerate}
		\item On r\'ep\`ete le lancer cinq fois successivement, de mani\`ere ind\'ependante.
			On consid\`ere comme succ\`es l'obtention d'une face dont le nombre inscrit se termine par $0$,
			et par cons\'equent comme \'echec tout autre \'eventualit\'e.
			\`A chaque succ\'es on inscrit le chiffre $1$ sur un tableau et \`a chaque \'echec on inscrit le chiffre $0$.
			On obtient de cette mani\`ere un nombre de cinq chiffres en base deux.
			(Exemple : \'echec--succ\`es--succ\`es--\'echec--succ\`es correspond \`a $\base{01101}{2}$.)
			On s'int\'eresse alors au nombre ainsi obtenu :
			calculer la probabilit\'e de chacun des \'ev\`enements suivants :
			\begin{description}
				\item[$A$ :] «~On obtient un nombre pair.~»
				\item[$B$ :] «~Le nombre obtenu contient exactement trois fois le chiffre $1$.~»
				\item[$C$ :] «~Le nombre \'ecrit sur le tableau est un multiple de $5$.~»
				\item[$D$ :] «~On obtient un nombre premier.~»
				(Indice : une suite contenant $k$ succ\`es a pour probabilit\'e
				$\left(\frac{1}{3}\right)^k\left(\frac{2}{3}\right)^{5-k}$.)
			\end{description}
	\end{enumerate}
\end{exo}


\section{Analyse}
\begin{exo}
Soit $f$ la fonction d\'efinie sur
$\mathcal{D}_f = [0 ; 1[ \cup ]1 ; +\infty[$ par
\[
  f(x)=
  \begin{cases}
    \dfrac{x}{\ln x} & \text{si } x \in ]0 ; 1[ \cup ]1 ; +\infty[, \\
    0                 & \text{si } x=0.
  \end{cases}
\]
Soit $(\mathscr{C})$ sa courbe repr\'esentative dans le plan muni d'un rep\`ere
orthonorm\'e \Oij{} d'unit\'e 1~cm.
\begin{enumerate}
\item \begin{enumerate}
      \item Montrer que $f$ est continue et d\'erivable au point $x_0 = 0$.
      \item Calculer les limites de $f$ aux bornes de $\mathcal{D}_f$.
      \end{enumerate}
\item \begin{enumerate}
      \item Calculer la fonction d\'eriv\'ee de $f$ et \'etudier son signe sur $\mathcal{D}_f$.
	      \item Dresser le tableau de variation de $f$.
	      \item \'Etudier les branches infinies de $f$ et tracer $(\mathscr{C})$ ainsi que 
            la droite $(T)$, tangente \`a $(\mathscr{C})$ au point d'abscisse $e^{-1}$ dans le rep\`ere \Oij. 
      \end{enumerate}
\item Soit $g$ la fonction d\'efinie sur $\mathcal{D}_g = ]1 ; +\infty[$ par
      $g(x) = \frac{x + \ln x}{\ln x}$.
      On note $(\Gamma)$ sa courbe repr\'esentative dans le rep\`ere \Oij.
      \begin{enumerate}
      \item \'Etudier la position relative de $(\Gamma)$ et $(\mathscr{C})$.
      \item Dresser le tableau de variation de $g$ et tracer $(\Gamma)$ dans le m\^eme rep\`ere que $(\mathscr{C})$.
      \item Calculer, en $\text{cm}^2$, l'aire du domaine plan limit\'e par les deux courbes $(\mathscr{C})$ et $(\Gamma)$ 
            et les droites d'\'equations respectives $x = 2$ et $x = e$.
      \end{enumerate}
\end{enumerate}
\end{exo}
\begin{exo}
On consid\`ere la fonction $f$ d\'efinie sur $I = [ 0 \; ; \ + \infty[$ par $f(x) = \frac{\expo{x} - 1}{\expo{x} - x}$
et on note $(\mathscr{C})$ sa courbe repr\'esentative dans un rep\`ere orthonorm\'e
\Oij{} d'unit\'e 2~cm.
\begin{enumerate}
\item Soit $h$ la fonction d\'efinie sur $I$ par : $h(x) = 2 \; \expo{x} - x \; \expo{x} - 1$.
      \begin{enumerate}
	      \item \'Etudier les variations de $h$.
      \item Montrer que l'\'equation $h(x) = 0$ admet une solution unique $\alpha$ dans $I$. 
            V\'erifier que ${\alpha \in \; \left] \frac{3}{2} \; ; \; 2 \right[}$.
      \item En d\'eduire le signe de $h(x)$ suivant les valeurs de $x$.
      \end{enumerate}
\item \begin{enumerate}
      \item D\'eterminer la limite de $f(x)$ quand $x$ tend vers $+\infty$.
            Interpr\'eter le r\'esultat graphiquement.
      \item \'Etudier la position de $(\mathscr{C})$ par rapport \`a la droite $(\mathcal{D})$ d'\'equation $y = 1$.
            (On admet que $\expo{x} > x $ pour tout $x \in I$)
      \end{enumerate}
\item \begin{enumerate}
      \item Montrer que pour tout $x \in I$, $f'(x) = \frac{h(x)}{\left(\expo{x} - x \right)^2}$.
      \item Montrer que $f(\alpha) = \frac{1}{\alpha - 1}$.
            Dresser le tableau des variations de $f$.
      \item Tracer $(\mathscr{C})$ et $(\mathcal{D})$ dans le m\^eme rep\`ere \Oij . 
            (On prendra $ \alpha \simeq 1,84$ et $f(\alpha) \simeq 1,20 $)
      \end{enumerate}
\item D\'eterminer la primitive $F$ de $f$ qui prend la valeur $e$ quand $x = 0$.
\end{enumerate}
\end{exo}
\begin{exo}
Soit $f$ la fonction d\'efinie sur $\mathbb{R}$ par $f(x) = x - 1 + 2 \; \expo{-\frac{x}{2}}$.
On note $(\mathscr{C})$ sa courbe repr\'esentative dans le plan rapport\'e \`a un
rep\`ere orthonorm\'e \Oij{} d'unit\'e 1~cm.
\begin{enumerate}
\item \begin{enumerate}
      \item \'Etudier les variations de $f$.
      \item Donner l'\'equation de l'asymptote oblique $(\mathcal{D})$ de $(\mathscr{C})$ et 
      pr\'eciser la position de $(\mathscr{C})$ par rapport \`a $(\mathcal{D})$.
      \item Tracer $(\mathscr{C})$ et $(\mathcal{D})$.
      \end{enumerate}
\item D\'eterminer la primitive de $f$ qui s'annule en $0$.
\item Soit $g$ la restriction de $f$ \`a l'intervalle $I = [0 \; ; \; +\infty[$.
      \begin{enumerate}
      \item Montrer que $g$ r\'ealise une bijection de $I$ vers un intervalle $J$ que l'on d\'eterminera.
      \item Calculer $g(1)$ et en d\'eduire $g^{-1} \left( \frac{2}{\sqrt{e}} \right) $.
      \end{enumerate}
\item Tracer dans le m\^eme rep\`ere que $(\mathscr{C})$ la courbe repr\'esentative de $g^{-1}$.
\end{enumerate}
\end{exo}
\begin{exo}
Soit $f$ la fonction d\'efinie sur $\mathbb{R}^*_+$ par $f(x) = \frac{1 + \ln (x)}{x}$.
On note $(\mathscr{C})$ sa courbe repr\'esentative dans un rep\`ere orthonorm\'e
\Oij{} d'unit\'e 2~cm.
\begin{description}
\item[A.] \begin{enumerate}
      \item Calculer les limites de $f$ aux bornes de son domaine de d\'efinition.
      \item Dresser le tableau de variations complet de $f$.
      \item D\'eterminer les coordonn\'ees du point $I$, intersection de $(\mathscr{C})$ avec l'axe des abscisses.
            \'Ecrire l'\'equation de la droite $(T)$ tangente \`a $(\mathscr{C})$ au point $I$.
      \item Tracer $(\mathscr{C})$ et $(T)$ dans le rep\`ere \Oij.
      \end{enumerate}
\item[B.] \begin{enumerate}
      \item Trouver la primitive $F$ de $f$ qui s'annule en $1$.
      \item D\'eterminer les limites de $F$ aux bornes de son domaine de d\'efinition puis dresser son tableau de variation.
      \item D\'eterminer les coordonn\'ees des points d'intersection de la courbe repr\'esentative de $F$ avec l'axe des abscisses.
      \item Repr\'esenter $F$ graphiquement. 
      \end{enumerate}
\end{description}
\end{exo}
\begin{exo}
      \begin{description}
            \item[I.] On consid\`ere la fonction $g$ telle que $g(x) = (2 x + 3) \expo{-2x} + 1$.
                  \begin{enumerate}
                        \item Dresser le tableau de variation de $g$.                                                                  
                        \item Montrer qu'il existe un r\'eel unique $\alpha$ dans $] -2 ; - \frac{3}{2} [$ 
                              tel que $g(\alpha) = 0$.      
                        \item En d\'eduire le signe de $g(x)$.                                                                         
                  \end{enumerate}
            \item[II.] On donne la fonction $f$ telle que $f(x) = (x+2) \expo{-2x} + 1 - x$.
                       On note $(\mathscr{C})$ sa courbe repr\'esentative dans un rep\`ere orthonorm\'e d'unit\'e 1~cm.
                  \begin{enumerate}
                        \item Exprimer $f'(x)$ en fonction de $g(x)$ o\`u $f'$ est la fonction d\'eriv\'ee de $f$.                    
                        \item En d\'eduire le tableau de variation de $f$.                                                             
                        \item \begin{enumerate}
                                    \item \'Etudier les branches infinies de $(\mathscr{C})$.                                                               
                                    \item Pr\'eciser la position de $(\mathscr{C})$ par rapport \`a son asymptote oblique $(\mathscr{D})$.                
                                    \item Montrer que l'\'equation $f(x) = 0$ admet une solution unique $a$ dans $]1 ; 2[$.                  
                                    \item Tracer $(\mathscr{C})$.                                                                                         
                              \end{enumerate}
                        \item Calculer l'aire du domaine plan limit\'e par $(\mathscr{C})$, 
                              l'axe des abscisses et les droites d'\'equations $x=0$ et $x=1$.                                         
                        \item On note $h$ la restriction de $f$ \`a l'intervalle $[0 ; +\infty[$.
                              \begin{enumerate}
                                    \item Montrer que $h$ r\'ealise une bijection de $[0 ; +\infty[$ 
                                          vers un intervalle $J$ que l'on pr\'ecisera.                                                       
                                    \item Calculer $(h^{-1})'(0)$.                                                                          
                                    \item Tracer la courbe $(\mathscr{C}')$, repr\'esentative de $h^{-1}$,
                                          dans le m\^eme rep\`ere que $(\mathscr{C})$.                
                              \end{enumerate}
                  \end{enumerate}
      \end{description}
\end{exo}
%%%%%%%%%%%%%%%%%%%%%%%%%%%%%%%%%%%%%%%%%%%%%%%%%%%%%%%%%%%%%%%%%%%%%%%%%%%%%%%%%%%
\begin{exo}
      On consid\`ere la fonction $f$ d\'efinie sur l'intervalle $[0 ; +\infty[$ par :
      \[
      f(x) = 
            \left\lbrace
                  \begin{array}{lcl}
                        x \ln \left( 1 + \frac{1}{x^2} \right) & \text{ si } & x \neq 0 \\
                        0                                      & \text{ si } & x = 0 \;.
                  \end{array}
            \right.
      \]
      On note $(\mathscr{C})$ la courbe repr\'esentative de $f$
      dans un rep\`ere orthonormal $(O ; \vec{\imath} , \vec{\jmath})$ (unit\'e graphique : 5~cm).
      \begin{description}
            \item[Partie A :]
                  On consid\`ere la fonction $g$ d\'efinie sur l'intervalle $]0 ; +\infty[$ par 
                  $g(x) = \ln \left( 1 + \frac{1}{x^2} \right) - \frac{2}{x^2 + 1}$.
                  \begin{enumerate}
                        \item \begin{enumerate}
                                    \item Calculer la d\'eriv\'ee $g'$ de $g$.
                                          Montrer que pour tout $x \in ]0 ; +\infty[$, 
                                          $g'(x) = \frac{2 (x^2 -1)}{x (x^2 + 1)^2}$.
                                    \item \'Etudier le signe de $g'(x)$ suivant les valeurs de $x$.
                              \end{enumerate}
                        \item D\'eterminer les limites de $g$ aux bornes de son domaine de d\'efinition.
                        \item \begin{enumerate}
                                    \item Dresser le tableau de variation de $g$.
                                    \item En d\'eduire qu'il existe un nombre r\'eel $\alpha > 0$ unique tel que 
                                          $g(\alpha) = 0$.
                                          V\'erifier que $0,5 < \alpha < 0,51$.
                              \end{enumerate}
                        \item D\'eduire des questions pr\'ec\'edentes le signe de $g(x)$ sur l'intervalle $]0 ; +\infty[$.
                  \end{enumerate}
            \item[Partie B :]
                  \begin{enumerate}
                        \item Montrer que pour tout $x\in ]0 ; +\infty[$ on a $f'(x) = g(x)$.
                              En d\'eduire les variations de $f$ sur $]0 ; +\infty[$.
                        \item \begin{enumerate}
                                    \item Calculer la limite quand $x$ tend vers $+\infty$ de $x f(x)$.
                                          (On pourra poser $X = \frac{1}{x^2}$.)
                                    \item En d\'eduire la limite de $f(x)$ quand $x$ tend vers $+\infty$.
                              \end{enumerate}
                        \item \begin{enumerate}
                                    \item D\'eterminer la limite de $f$ en $0$.
                                          (On pourra \'ecrire $f(x)$ sous la forme $ x \ln (x^2+1) - 2 x \ln x$.)
                                    \item \'Etudier la d\'erivabilit\'e de $f$ en $0$.
                                          Pr\'eciser la tangente \`a la courbe $(\mathscr{C})$ en $0$.
                              \end{enumerate}
                        \item \begin{enumerate}
                                    \item Prouver que, pour tout $x \in [ 0,5 ; \alpha ]$,
                                          $ 0 \leq f'(x) \leq f'(0,5)$.
                                    \item En d\'eduire que, pour tout \'el\'ement $x$ de $[0,5 ; \alpha]$, 
                                          $0 \leq f(\alpha) - f(0,5) \leq (\alpha - 0,5) f'(0,5)$,
                                          puis que 
                                          $0 \leq f(\alpha) - f(0,5) \leq \frac{1}{100} f'(0,5)$.
                                    \item En d\'eduire une valeur approch\'ee de $f(\alpha)$ \`a $10^{-3}$ pr\`es.
                              \end{enumerate}
                        \item Dresser le tableau de variation de $f$ puis tracer la courbe $(\mathscr{C})$.
                  \end{enumerate}
            \item[Partie C :]
                  Soit $\lambda$ un nombre r\'eel appartenant \`a l'intervalle $]0 ; 1]$.
                  \begin{enumerate}
                        \item En utilisant une int\'egration par parties, calculer l'int\'egrale
                              $J_\lambda = \int_{\lambda}^{1} f(x) \dx$.
                              Donner une interpr\'etation g\'eom\'etrique de cette int\'egrale.
                        \item Calculer $\displaystyle{\lim_{\lambda \rightarrow 0} J_\lambda}$.
                              On admet que cette limite est l'aire de la partie du plan constitu\'ee
                              des points dont les coordonn\'ees $(x ; y)$ v\'erifient : 
                              $
                              \left\lbrace
                                    \begin{array}{l}
                                          0 \leq x \leq 1    \\
                                          0 \leq y \leq f(x) \;.
                                    \end{array}
                              \right.
                              $
                              En d\'eduire la valeur de cette aire exprim\'ee en $\mathrm{cm}^2$.
                  \end{enumerate}
      \end{description}
\end{exo}
\begin{exo} 
      \begin{description}
      \item[Partie A :]
            Pour tout entier naturel $n>0$, soit la fonction num\'erique $f_n$ de la variable r\'eelle $x$
            d\'efinie sur $\mathbb{R} \setminus \{ -1 \}$ par :
            \[
                  f_n(x) = \frac{\expo{-x}}{(1+x)^n}.
            \]
            On note $(\mathscr{C}_n)$ la courbe repr\'esentative de $f_n$ dans un rep\`ere orthonorm\'e
            $(O ; \vec{\imath} , \vec{\jmath} )$ d'unit\'e 1~cm.
            \begin{enumerate}
                  \item \begin{enumerate}
                              \item Calculer les limites de $f_n$ aux bornes de son domaine de d\'efinition
                                    (on distinguera le cas o\`u $n$ est pair ou impair).
                              \item Calculer $f'_n (x)$ en fonction de $x$ et de $n$.
                              \item Dresser le tableau de variation de $f_n$ selon la parit\'e de $n$.
                        \end{enumerate}
                  \item Montrer que toutes les courbes $(\mathscr{C}_n)$ passent par un point 
                        fixe dont on pr\'ecisera les coordonn\'ees.
                  \item Calculer $\displaystyle{\lim_{x\rightarrow-\infty} \frac{f_n(x)}{x}}$
                        (distinguer le cas o\`u $n$ est pair ou impair).
                  \item \begin{enumerate}
                              \item \'Etudier la position relative des courbes $(\mathscr{C}_1)$ et $(\mathscr{C}_2)$.
                              \item Tracer $(\mathscr{C}_1)$ et $(\mathscr{C}_2)$ dans le m\^eme rep\`ere.
                        \end{enumerate}
            \end{enumerate}

      \item[Partie B :]
            Pour tout entier naturel $n > 0$, soit $\displaystyle{I_n = \int_{0}^{1} f_n (x) \dx }$.
            \begin{enumerate}
                  \item Donner l'expression de $f_n'(x)$ en fonction de $f_n (x)$ et de $f_{n+1} (x)$.
                  \item \begin{enumerate}
                              \item Montrer que la suite $(I_n)$ est d\'ecroissante.
                              \item En d\'eduire que la suite $(I_n)$ est convergente.
                        \end{enumerate}
                  \item \begin{enumerate}
                              \item D\'emontrer que, pour tout entier naturel $n > 0$ et tout $x \in [0 ; 1]$, on a
                                    \[
                                      \frac{\expo{-1}}{(1+x)^n}
                                      \leq f_n(x)
                                      \leq \frac{1}{(1+x)^n}.
                                    \]
                              \item En d\'eduire que, pour tout $n > 1$, on a
                                    \[
                                      \begin{aligned}
                                        \frac{\expo{-1}}{n-1}
                                        \left(1-\frac{1}{2^{n-1}}\right)
                                        &\leq I_n, \\
                                        I_n
                                        &\leq \frac{1}{n-1}
                                        \left(1-\frac{1}{2^{n-1}}\right).
                                      \end{aligned}
                                    \]
                              \item Calculer $\displaystyle{\lim_{n\rightarrow+\infty} I_n}$.
                        \end{enumerate}
                  \item \begin{enumerate}
                              \item En utilisant la question 1 de la Partie B, montrer que 
                                    $I_n + n I_{n+1} = 1 - \frac{\expo{-1}}{2^n}$.
                              \item En d\'eduire $\displaystyle{\lim_{n\rightarrow+\infty} n I_{n+1}}$.
                        \end{enumerate}
            \end{enumerate}
      \end{description}
\end{exo}
\begin{exo}
      \begin{description}
      \item[Partie A :]
      Soit $f$ la fonction d\'efinie sur $[0 ; +\infty[$ par
      $\displaystyle f(x) = \frac{\expo{-x^2}}{1+x}$.
      On note $(\mathscr{C})$ sa courbe repr\'esentative dans un rep\`ere
      orthonorm\'e d'unit\'e 2~cm.
      \begin{enumerate}
      \item \'Etudier les variations de $f$ et dresser son tableau de variations.
      \item Tracer la courbe $(\mathscr{C})$.
      \item Dans cette question, on se propose d'encadrer l'int\'egrale
            \[
              K = \int_0^1 \frac{\expo{-x^2}}{1+x} \dx.
            \]
            \begin{enumerate}
            \item Interpr\'eter g\'eom\'etriquement l'int\'egrale $K$.
            \item En \'etudiant les variations des fonctions 
                  $g : x \mapsto \expo{-x} + x - 1$ et 
                  $h : x \mapsto 1 - x + \frac{x^2}{2} - \expo{-x}$
                  sur l'intervalle $[0 , 1]$, d\'emontrer que pour tout $x $ de $[0 , 1]$,
                  \begin{equation} \label{eq:first}
                  1 - x \leq \expo{-x} \leq 1 - x + \frac{x^2}{2} \; .
                  \end{equation}
            \item En d\'eduire un encadrement de $\expo{-x^2}$ pour $x$ \'el\'ement de $[0 , 1]$ 
                  puis montrer que pour tout $x$ de 
                  $[0 , 1]$ : 
                  \begin{equation} \label{eq:second}
                  1-x \leq \frac{\expo{-x^2}}{1 + x} \leq 1 - x + \frac{x^4}{2 \cdot (1 + x)} \; .
                  \end{equation} 
            \end{enumerate}
            \item Montrer que, pour tout $x$ de $[0 , 1]$,
                  \[
                    \frac{x^4}{1 + x} = x^3 - x^2 + x - 1 + \frac{1}{1 + x}.
                  \]
            \item D\'eduire alors de \eqref{eq:second} que
                  \[
                    \frac{1}{2} \leq K \leq \frac{5}{24} + \frac{\ln 2}{2}.
                  \]
      \end{enumerate}
      \item[Partie B :]
      Dans cette partie, 
      on \'etudie la convergence de la suite $(u_n)_{n \in \mathbb{N}^*}$ d\'efinie, pour tout entier naturel non nul $n$, par : 
      \[ 
        u_n = \int_0^1 \frac{\expo{-x^2}}{(1 + x)^n} \dx.
      \]
      \begin{enumerate}
      \item \'Etudier la variation de la suite $(u_n)$.
      \item En d\'eduire que la suite $(u_n)$ est convergente.     
      \item Montrer que, pour tout $n \geq 2$, on a
            \[
              \begin{aligned}
                \frac{1}{e(n-1)}\left(1-\frac{1}{2^{n-1}}\right)
                &\leq u_n, \\
                u_n
                &\leq \frac{1}{n-1}\left(1-\frac{1}{2^{n-1}}\right).
              \end{aligned}
            \]
      \item En d\'eduire la limite de la suite $(u_n)$.      
      \end{enumerate}
      \end{description}
\end{exo}
\begin{exo}
      On d\'esigne par $f$ la fonction num\'erique d\'efinie sur $]0 ; +\infty[$ par $f(x) = \frac{\ln x}{\sqrt{x}}$.
      On note $(\mathscr{C})$ la courbe repr\'esentative de $f$ dans un rep\`ere
      $(O;\vec{\imath},\vec{\jmath})$ d'unit\'e 1~cm.
      \begin{description}
            \item[I. \'Etude de la fonction $f$.] ~\\[-1.5\baselineskip]
            \begin{enumerate}
                  \item Calculer les limites de $f$ aux bornes de son domaine de d\'efinition
                        puis dresser le tableau de variation de $f$.
                  \item On d\'esigne par $(T)$ la tangente \`a $(\mathscr{C})$ au point d'abscisse 1.
                        \begin{enumerate}
                              \item D\'eterminer une \'equation de $(T)$.
                              \item On d\'esigne par $g$ la fonction d\'efinie sur $]0;+\infty[$ par 
                                    $g(x) = (x-1) - f(x)$.
                                    Montrer que 
                                    $g'(x) = \displaystyle{
                                          \frac{1}{2x\sqrt{x}}
                                                \left[
                                                      \ln x + 2 \left( x \sqrt{x} - 1 \right)
                                                \right]
                                          }$.
                                    Calculer $g'(1)$ et \'etudier le signe de $g'$ sur chacun des intervalles
                                    $]0;1[$ et $]1;+\infty[$.
                              \item Calculer $g(1)$ et, \`a l'aide du sens de variation de $g$, \'etudier le signe de $g(x)$.
                                    En d\'eduire la position de $(\mathscr{C})$ par rapport \`a $(T)$.
                        \end{enumerate}
                  \item Tracer la courbe $(\mathscr{C})$ ainsi que ses asymptotes.
            \end{enumerate}
            \item[II. \'Etude de suites.] ~\\
            Pour tout entier naturel $n$ sup\'erieur ou \'egal \`a $8$, on pose
            $u_n = \sum_{k=8}^n f(k) = f(8) + f(9) + \cdots + f(n)$.
            \begin{enumerate}
                  \item \label{quest:pbm1}
                         \begin{enumerate}
                              \item Soit $k$ un entier sup\'erieur ou \'egal \`a $8$. \label{quest:pbm1a}
                                    D\'emontrer que : $f(k+1) \leq \int_{k}^{k+1} f(t) \dt \leq f(k)$.
                              \item En d\'eduire que, pour tout entier $n \geq 8$, 
                                    $u_{n+1} - f(8) \leq \int_{8}^{n+1} f(t) \dt \leq u_n$.
                              \item \`A l'aide d'une int\'egration par parties, calculer $I_n = \int_{8}^{n+1} f(t) \dt$.
                              \item D\'eduire des questions pr\'ec\'edentes la limite de $(u_n)$.
                        \end{enumerate}
                  \item Pour tout entier naturel $n \geq 8$, on pose 
                        $v_n = u_n - \int_{8}^{n+1} f(t) \dt$.
                        \begin{enumerate}
                              \item \'Etudier le sens de variation de $(u_n)$ et en d\'eduire que 
                                    $u_n - f(8) \leq \int_{8}^{n+1} f(t) \dt \leq u_n$
                                    puis montrer que la suite $(v_n)$ est born\'ee.
                              \item En utilisant la question \ref{quest:pbm1}~\ref{quest:pbm1a} 
                                    \'etudier le sens de variation de $(v_n)$.
                              \item Justifier que la suite $(v_n)$ est convergente et d\'emontrer que sa limite
                                    $\ell$ v\'erifie : $0 \leq \ell \leq 0,74$.
                        \end{enumerate}
            \end{enumerate}
      \end{description}
\end{exo}
%
\begin{exo}
      \begin{description}
            \item[Partie A :]
                  On consid\`ere la fonction $g$ d\'efinie sur l'intervalle
                  $I = \left[ -\frac{1}{2} ; +\infty \right[$ par 
                  $g(x) = \ln (1+x) - x + \frac{x^2}{2} - \frac{x^3}{3}$.
                  \begin{enumerate}
                        \item \begin{enumerate}
                                    \item D\'emontrer que, pour tout $t$ de $I$, on a : $g'(t) = - \frac{t^3}{1 + t}$.
                                    \item Montrer alors que, pour tout $t$ de $I$, on a : 
                                          $\left| g'(t) \right| \leq 2 |t|^3$.
                              \end{enumerate}
                        \item D\'emontrer que, pour tout $x$ de $I$, on a :
                              $|g(x)| \leq \frac{x^4}{2}$.
                  \end{enumerate}
            \item[Partie B :]
                  Soit $f$ la fonction d\'efinie sur $]-1 ; +\infty[$ par 
                  \[
                        f(x) = 
                        \left\lbrace
                              \begin{array}{ll}
                                    \frac{x-\ln(1+x)}{x^2} & \text{ si } x \neq 0 \\
                                    \frac{1}{2} & \text{ si } x = 0.
                              \end{array}
                        \right.
                  \]
                  On note $(\mathscr{C})$ la courbe repr\'esentative de $f$ dans un rep\`ere orthogonal $(O ; \vec{\imath} , \vec{\jmath})$ d'unit\'e 2~cm.
                  \begin{enumerate}
                        \item En utilisant les encadrements obtenus dans la partie~A, d\'emontrer que la fonction 
                              $f$ est d\'erivable en $0$ et donner une \'equation de la tangente \`a $(\mathscr{C})$ au point $A(0 ; \frac{1}{2})$.
                              Montrer que $f$ est continue en $0$.
                        \item Soit $h$ la fonction d\'efinie sur $]-1 ; +\infty[$ par : 
                              $h(x) = \frac{-x^2-2x}{1+x}+2\ln(1+x)$.
                              \begin{enumerate}
                                    \item \'Etudier le sens de variation de $h$ (on ne demande pas les limites aux bornes).
                                    \item Pr\'eciser $h(0)$. En d\'eduire le signe de $h(x)$ sur $]-1 ; +\infty[$.
                              \end{enumerate}
                        \item Calculer $f'(x)$ pour $x \in ]-1 ; +\infty[ \setminus \{0\}$ et exprimer $f'(x)$ en fonction de $x$ et de $h(x)$.
                              En d\'eduire la variation de $f$.
                        \item \'Etudier les limites de $f$ en $-1$ et en $+\infty$.
                        \item Construire la courbe $(\mathscr{C})$ en pr\'ecisant ses asymptotes.
                  \end{enumerate}
            \item[Partie C :]
                  \begin{enumerate}
                        \item \begin{enumerate}
                                    \item D\'emontrer que, pour tout r\'eel $t$ positif ou nul, on a : $-t^2 \leq h'(t) \leq 0$.
                                    \item Pour tout r\'eel $x$ positif ou nul, 
                                    en d\'eduire par int\'egration un encadrement de $h(x)$, puis d\'emontrer que $|f'(x)| \leq \frac{1}{3}$.
                              \end{enumerate}
                        \item On d\'efinit la fonction $\varphi$ sur $]0 ; +\infty[$ par : $\varphi (x) = f(x) - x$.
                              De l'\'etude de la variation de $\varphi$, d\'eduire que l'\'equation $f(x) = x$ admet une solution r\'eelle strictement positive qu'on notera $\alpha$.
                              V\'erifier que $\alpha < 1$.
                  \end{enumerate}
            \item[Partie D :]
                  \begin{enumerate}
                        \item On pose $u_0 = 1$ et $u_{n+1} = f(u_n)$ pour tout entier naturel $n$.
                              Montrer, en utilisant la variation de $f$, que cette suite est bien d\'efinie et que,
                              pour tout entier naturel $n$, on a : $0 \leq u_n \leq 1$.
                        \item \begin{enumerate}
                                    \item D\'emontrer que, pour tout entier naturel $n$, on a :
                                          $| u_{n+1} - \alpha | \leq \frac{1}{3} |u_n - \alpha|$.
                                          Montrer ensuite que : 
                                          $|u_n - \alpha | \leq \frac{1}{3^n}$.
                                    \item Prouver que la suite $(u_n)$ converge et pr\'eciser sa limite.
                                    \item Donner une valeur approch\'ee de $\alpha$ \`a $10^{-2}$ pr\`es.
                              \end{enumerate}
                  \end{enumerate}
      \end{description}
\end{exo}
%
\begin{exo}
      \begin{description}
            \item[Partie A :]
                  On consid\`ere la fonction $f$ d\'efinie sur $[0 ; +\infty[$ par : 
                  \[
                        f(x) = 
                              \left\lbrace
                              \begin{array}{ll}
                                    x \; \ln \left( \frac{x+2}{x} \right) + \frac{x}{4} + \frac{1}{2} & \text{ si } x > 0 \\
                                    \frac{1}{2} & \text{ si } x = 0.
                              \end{array}
                              \right.
                  \]
                  On note $(\mathscr{C})$ sa courbe repr\'esentative dans un rep\`ere orthonormal $(O ; \vec{\imath} , \vec{\jmath})$ d'unit\'e graphique 2~cm.
                  \begin{enumerate}
                        \item Soit $g$ la fonction d\'efinie sur $]0 ; +\infty[$ par 
                              $g(x) = \ln (x+2) - \ln (x) - \frac{2}{x+2} + \frac{1}{4}$.
                              \begin{enumerate}
                                    \item \begin{enumerate}
                                                \item \'Etudier la variation de $g$.
                                                \item D\'eterminer $\displaystyle{\lim_{x\rightarrow +\infty} g(x)}$.
                                                \item En d\'eduire le signe de $g$ sur $]0 ; +\infty[$.
                                          \end{enumerate}
                                    \item Montrer que pour tout $x \in [ 2 ; 3 ]$ on a $g(x) < \frac{1}{2}$.
                              \end{enumerate}
                        \item \begin{enumerate}
                                    \item D\'eterminer la limite quand $x$ tend vers $0$ par valeur sup\'erieure de
                                          $x \; \ln \left( \frac{x+2}{x} \right)$ (on pourra poser $t = \frac{1}{x}$)
                                          et d\'emontrer que $f$ est continue en $0$.
                                    \item Calculer $\displaystyle{\lim_{x\rightarrow 0^+} \frac{f(x) - f(0)}{x-0}}$.
                                          Interpr\'eter ce r\'esultat.
                                    \item \'Etudier les variations de $f$.
                                    \item \begin{enumerate}
                                                \item D\'emontrer que 
                                                      $\displaystyle{
                                                      \lim_{x\rightarrow+\infty} x \; \ln \left( \frac{x+2}{x} \right) = 2
                                                      }$.
                                                \item En d\'eduire la limite de $f(x)$ quand $x$ tend vers $+\infty$.
                                                \item Montrer que la droite $\Delta$ d'\'equation 
                                                      $ y = \frac{x}{4} + \frac{5}{2}$
                                                      est asymptote \`a la courbe $(\mathscr{C})$ au voisinage de $+\infty$.
                                          \end{enumerate}
                                    \item Tracer $(\mathscr{C})$, la droite $\Delta$ ainsi que la droite $(y=x)$ 
                                          dans le m\^eme rep\`ere.
                              \end{enumerate}
                  \end{enumerate}
            \item[Partie B :]
                  Dans cette partie, on d\'esigne par $I$ l'intervalle $[2;3]$.
                  \begin{enumerate}
                        \item \begin{enumerate}
                                    \item Soit la fonction $h$ d\'efinie sur $I$ par $h(x) = f(x) - x$.
                                          Montrer que pour tout $x \in I$, $h'(x) < 0$
                                          (remarquer que $h'(x) = g(x) - 1$).
                                    \item En d\'eduire la variation de $h$ et montrer que l'\'equation $h(x) = 0$
                                          admet une solution unique $\alpha$ dans $I$.
                              \end{enumerate}
                        \item \begin{enumerate}
                                    \item Montrer que pour tout $x \in I$, $ 0 < f'(x) < \frac{1}{2}$.
                                    \item En d\'eduire que pour tout $x \in I$, 
                                          $|f(x) - \alpha | \leq \frac{1}{2} | x- \alpha|$.
                              \end{enumerate}
                  \end{enumerate}
            \item[Partie C :]
                  Soit $(u_n)$ la suite num\'erique d\'efinie pour tout entier naturel $n$ par :       
                  \[
                        \left\lbrace
                              \begin{array}{l}
                                    u_0 = \frac{9}{4}            \\
                                    u_{n+1} = f(u_n)             \;.
                              \end{array}
                        \right.
                  \]
                  \begin{enumerate}
                        \item Montrer que pour tout $n\in\mathbb{N}$, $u_n \in I$.
                        \item Montrer, en utilisant l'in\'egalit\'e des accroissements finis, que pour tout $n$, on a :
                              \[
                                    | u_{n+1} - \alpha | \leq \frac{1}{2} | u_n - \alpha |.
                              \]
                        \item Montrer que la suite $(u_n)$ converge vers $\alpha$.
                  \end{enumerate}
      \end{description}
\end{exo}
%
\begin{exo}
      On consid\`ere la fonction $f$ d\'efinie par $f(x) = \displaystyle{ 3 - x + \frac{\ln |x|}{x^3} }$.
      On note $(\mathscr{C})$ la courbe repr\'esentative de $f$ dans un rep\`ere 
      orthonormal $(O ; \vec{\imath} , \vec{\jmath} )$ d'unit\'e 2~cm.
      \begin{description}
            \item[Partie A :]
                  \begin{enumerate}
                        \item Soit $h$ la fonction d\'efinie sur $\mathbb{R}^*$ par : 
                              $h(x) = - x^4 + 1 - 3 \ln |x|$.
                              \begin{enumerate}
                                    \item Calculer $h(1)$ et $h(-1)$.
                                    \item Calculer $h'(x)$ et dresser le tableau de variation de $h$.
                                          En d\'eduire le signe de $h(x)$ suivant les valeurs de $x$.
                              \end{enumerate}
                        \item \begin{enumerate}
                                    \item D\'eterminer le domaine de d\'efinition $D_f$ de $f$.
                                    \item Calculer $f'(x)$ et montrer que $f'(x) = h(x) \cdot x^{-4}$.
                                    \item Dresser le tableau de variation de $f$.
                                    \item Montrer que l'\'equation $f(x) = 0$ admet exactement deux solutions 
                                          positives $\alpha$ et $\beta$ telles que
                                          $\alpha \in [\frac{1}{2} ; 1]$ et $\beta \in [3 ; 4]$.
                              \end{enumerate}
                        \item \begin{enumerate}
                                    \item Montrer que $(\mathscr{C})$ poss\`ede une asymptote oblique $\Delta$ 
                                          dont on pr\'ecisera l'\'equation.
                                    \item \'Etudier la position relative de $(\mathscr{C})$ par rapport \`a $\Delta$.
                                    \item D\'eterminer les coordonn\'ees du point $A$ d'abscisse positive
                                           tel que la tangente \`a $(\mathscr{C})$ au point $A$ est parall\`ele \`a $\Delta$.
                              \end{enumerate}
                        \item Tracer $(\mathscr{C})$ et $\Delta$.
                  \end{enumerate}
            \item[Partie B :]
                  On consid\`ere la fonction $g$ d\'efinie pour tout $x \geq 3$ par 
                  $g(x) = \displaystyle{ 3 + \frac{\ln x}{x^3} }$.
                  \begin{enumerate}
                        \item \begin{enumerate}
                                    \item Calculer $g'(x)$ et $g''(x)$.
                                          Dresser les tableaux de variation de $g'$ et de $g$.
                                    \item Montrer que $g\left( [ 3 ; 4 ] \right) \subset [ 3 ; 4 ]$.
                              \end{enumerate}
                        \item Tracer sur le m\^eme rep\`ere que $(\mathscr{C})$ la courbe 
                              repr\'esentative $(\Gamma)$ de $g$.
                        \item \begin{enumerate}
                                    \item Montrer que l'\'equation $g(x) = x$ admet une solution $\gamma \in [ 3 ; 4 ]$.
                                    \item Que peut-on dire de $\beta$ et de $\gamma$\,?
                              \end{enumerate}
                        \item \begin{enumerate}
                                    \item Montrer que $\forall x \in [ 3 ; 4 ]$, on a $g'(x) \leq 0$
                                          et que $|g'(x)| \leq \frac{3}{x^3}$ (indication : $\ln x \leq x$).
                                    \item En d\'eduire que $\forall x \in [ 3 ; 4 ]$, $|g'(x)| \leq \frac{1}{9}$
                                          (indication : $x \mapsto \frac{1}{x^3}$ est d\'ecroissante).
                                    \item Montrer que $\forall (a,b) \in [3;4] \times [3;4]$,
                                          $
                                                \left| g(b) - g(a) \right| 
                                                      \leq 
                                                \frac{1}{9} |b-a|
                                          $.
                              \end{enumerate}
                        \item On consid\`ere la suite $(u_n)$ d\'efinie par 
                              \[
                                    \left\lbrace
                                          \begin{array}{l}
                                                u_0 = 3 \\
                                                u_{n+1} =
                                                      \displaystyle{ 3+\frac{\ln (u_n)}{\left( u_n \right)^3} 
                                                                    }
                                                                    \quad \forall n\in\mathbb{N}
                                          \end{array}
                                    \right.
                              \]
                              \begin{enumerate}
                                    \item Montrer par r\'ecurrence que $\forall n \geq 0$, $3 \leq u_n \leq 4$.
                                    \item Montrer que $\forall n \geq 0$,
                                          $
                                          |u_{n+1} - \gamma | \leq \frac{1}{9} |u_n - \gamma |
                                          $
                                          et que 
                                          $
                                          |u_n - \gamma | \leq \left( \frac{1}{9} \right)^n.
                                          $
                                    \item Donner une valeur approch\'ee de $\gamma$ \`a $10^{-2}$ pr\`es.
                                    \item Calculer la limite de $u_n$ quand $n$ tend vers $+\infty$.
                              \end{enumerate}
                  \end{enumerate}
      \end{description}
\end{exo}
%
\begin{exo}
      \begin{description}
            \item[Partie A :]
                  On consid\`ere la fonction $f$ d\'efinie sur $\mathbb{R}$ par $f(x) = \ln \left( 1 + \expo{-2x} \right)$.
                  On d\'esigne par $(\mathscr{C})$ sa courbe repr\'esentative dans le plan rapport\'e \`a un rep\`ere 
                  orthonormal $(O ; \vec{\imath} , \vec{\jmath} )$ d'unit\'e 2~cm.
                  \begin{enumerate}
                        \item \'Etudier les variations de $f$ et dresser son tableau de variation.
                        \item \begin{enumerate}
                                    \item Montrer que pour tout $x$ r\'eel, 
                                          $f(x) = \ln \left( 1 + \expo{2x} \right) - 2x $.
                                    \item En d\'eduire que la courbe $(\mathscr{C})$ admet comme asymptote la droite 
                                          $(\mathscr{D})$ d'\'equation $y = -2x$.
                                    \item \'Etudier la position relative de $(\mathscr{C})$ et $(\mathscr{D})$.
                              \end{enumerate}
                        \item Tracer la droite $(\mathscr{D})$ et la courbe $(\mathscr{C})$.
                        \item \begin{enumerate}
                                    \item Montrer que $f$ r\'ealise une bijection 
                                          de $\mathbb{R}$ sur un intervalle $J$ \`a pr\'eciser.
                                    \item D\'eterminer l'expression explicite de $f^{-1} (x)$.
                              \end{enumerate}
                  \end{enumerate}
            \item[Partie B :]
                  Pour tout $x$ appartenant \`a l'intervalle $[ 0 ; +\infty[$, 
                  on pose $F(x) = \displaystyle{\int_{0}^{x} \ln \left( 1 + \expo{-2t} \right) \dt }$.
                  (On ne cherchera pas \`a calculer $F(x)$).
                  \begin{enumerate}
                        \item Pour $x > 0$, donner une interpr\'etation g\'eom\'etrique de $F(x)$.
                        \item \'Etudier le sens de variation de $F$ sur $[0 ; +\infty[$.
                        \item Soit $a$ un r\'eel strictement positif. \label{quest:log-inequality}
                              \begin{enumerate}
                                    \item Montrer que pour tout $t \in [1 ; 1 + a ]$, on a 
                                          $\frac{1}{1+a} \leq \frac{1}{t} \leq 1$.
                                    \item En appliquant le th\'eor\`eme des accroissements finis
                                          \`a la fonction logarithme n\'ep\'erien, \'etablir que : 
                                          \[
                                                \frac{a}{1+a} \leq \ln (1 + a ) \leq a .
                                          \]
                              \end{enumerate}
                        \item Soit $x$ un r\'eel strictement positif.
                              D\'eduire de la question~\ref{quest:log-inequality} que 
                              \[
                                \begin{aligned}
                                  \int_{0}^{x}\frac{\expo{-2t}}{1+\expo{-2t}}\dt
                                  &\leq F(x), \\
                                  F(x) &\leq \int_{0}^{x}\expo{-2t}\dt.
                                \end{aligned}
                              \]
                              Puis
                              \[
                                \begin{aligned}
                                  \frac{1}{2}\ln 2
                                  -\frac{1}{2}\ln\left(1+\expo{-2x}\right)
                                  &\leq F(x), \\
                                  F(x) &\leq \frac{1}{2}-\frac{1}{2}\expo{-2x}.
                                \end{aligned}
                              \]
                        \item On admet que la limite de $F(x)$ en $+\infty$ existe et est un nombre r\'eel $\ell$.
                              \'Etablir que $\frac{1}{2} \ln 2 \leq \ell \leq \frac{1}{2}$.
                        \item Pour tout entier naturel $n$, on pose 
                              $
                              u_n = \displaystyle{
                                          \int_{n}^{n+1} \ln \left( 1 + \expo{-2t} \right) \dt
                                    }
                              $.
                              \begin{enumerate}
                                    \item Montrer que pour tout entier naturel $n$, on a 
                                          $0 \leq u_n \leq \ln (1 + \expo{-2n})$
                                          (Utiliser le sens de variation de $f$).
                                    \item D\'eterminer la limite de la suite $(u_n)$.
                              \end{enumerate}
                        \item Pour tout entier naturel $n$, on pose $S_n = u_0 + u_1 + \cdots + u_n$.
                              \begin{enumerate}
                                    \item Exprimer $S_n$ en fonction de $F$ et de $n$.
                                    \item Montrer que la suite $(S_n)$ converge et pr\'eciser sa limite.
                              \end{enumerate}
                  \end{enumerate}
      \end{description}
\end{exo}
%%%%%%%%%%%%%%%%%%%%%%%%%%%%%%%%%%%%%%%%%%%%%%%%%%%%%%%%%%%%%%%%%%%%%%%%%%%%%%%%%%%

